% *********************************************************************
% © 2010–2021 Jeremy Sylvestre
%
% Permission is granted to copy, distribute and/or modify this document
% under the terms of the GNU Free Documentation License, Version 1.3 or
% any later version published by the Free Software Foundation; with no
% Invariant Sections, no Front-Cover Texts, and no Back-Cover Texts. A
% copy of the license is included in the appendix entitled “GNU Free
% Documentation License” that appears in the output document of this
% PreTeXt source code. All trademarks™ are the registered® marks of
% their respective owners.
%
% *********************************************************************

\begin{tikzpicture}[
	vertex/.style={ circle, draw, very thin, fill, inner sep=0pt, minimum size=4pt },
	vertexg/.style={ black!40!white, circle, draw, very thin, fill, inner sep=0pt, minimum size=4pt },
	vertexo/.style={ circle, draw, very thin, inner sep=0pt, minimum size=4pt }
]

% old command macros
% \vertex            TO  \node at (#2) (#3) [vertex#1] {};
% \lvertex           TO  \node at (#2) (#3) [vertex#1,label=#4] {};
% \vertexarrowto     TO  \node at (#2) (#3) [vertex#1] {};
%                        \foreach \p in {#4} \draw[-latex] (#3) to (\p);
% \lvertexarrowto    TO  \node at (#2) (#3) [vertex#1,label=#4] {};
%                        \foreach \p in {#5} \draw[-latex] (#3) to (\p);
% \vertexlinefrom    TO  \node at (#2) (#3) [vertex#1] {};
%                        \foreach \p in {#4} \draw (\p) to (#3);
% \vertexllinefrom   TO  \vertex[#1]{#2}{#3}
%                        \foreach \p in {#4} \draw (\p) to node[#6] {#5} (#3);
% \lvertexllinefrom  TO  \vertexllinefrom[#1,label=#4]{#2}{#3}{#5}{#6}{#7}
%
% \lvertexline     ALIAS \vertexlinefrom[#1,label=#4]{#2}{#3}{#5}
% \vertexarrowfrom   TO  \vertex[#1]{#2}{#3}
%                        \foreach \p in {#4} \draw[-latex] (\p) to (#3);
%
%
% ALL OF THESE HAVE DEFAULT OPTIONAL ARG = 0.25cm
% \loopsouth       TO  \draw (#2) arc (85:-265:#1) to (#2);
% \loopnorth       TO  \draw (#2) arc (-85:265:#1) to (#2);
% \loopwest        TO  \draw (#2) arc (5:355:#1) to (#2);
% \loopeast        TO  \draw (#2) arc (175:-175:#1) to (#2);
% \arrowloopsouth  TO  \draw[-latex] (#2) arc (85:-225:#1) to (#2);
% \arrowloopnorth  TO  \draw[-latex] (#2) arc (-85:225:#1) to (#2);
% \arrowloopwest   TO  \draw[-latex] (#2) arc (5:315:#1) to (#2);
% \arrowloopeast   TO  \draw[-latex] (#2) arc (175:-135:#1) to (#2);



\end{tikzpicture}
